\section{Related work}\label{sec:related}
\paragraph{Benchmark prediction from few items.}
Anchor Points~\citep{vivek2024anchor} cluster items by the response vectors of source models and evaluate a new model on one representative per cluster, weighting by cluster size. tinyBenchmarks~\citep{polo2024tiny} combine such anchors with a multidimensional item-response-theory (IRT) model to impute the unseen items; metabench~\citep{kipnis2025metabench} selects items by IRT information and predicts with generalised additive models. \citet{pacchiardi2024100} train assessor models on a few reference instances, and \citet{li2025active} learn an adaptive item-acquisition policy. \citet{perlitz2024efficient} study the same cost problem from a benchmark-reduction view, and \citet{rodriguez2021evaluation} use IRT to show that items differ widely in informativeness. \citet{zhang2025misses} compare eleven such methods on 19 benchmarks and report that most of them interpolate well among similar models but struggle when the target differs from the models seen before, whereas random sampling and an AIPW estimator stay reliable (we know this work only through its description in the task's starter list and the README of its data repository; see Appendix~\ref{app:repro}). Our experiments use the item-level data released with that work, but our protocols, estimators and designs are implemented independently.

\paragraph{Survey sampling and prediction-powered inference.}
Estimating a finite-population mean from a sample is a classical problem. Horvitz--Thompson weighting~\citep{horvitz1952generalization} gives design-unbiased estimates under unequal-probability sampling, Neyman allocation~\citep{neyman1934two} minimises variance given stratum variances, and model-assisted (GREG/difference) estimators~\citep{sarndal1992model} add a design-weighted residual correction to model predictions, so that the estimate remains unbiased however wrong the model is. The same construction underlies doubly-robust/AIPW estimation~\citep{robins1994estimation} and prediction-powered inference~\citep{angelopoulos2023ppi}. We treat benchmark prediction as a finite-population problem where the items are the population, the (unknown) target model's correctness is the variable of interest, and source models provide the auxiliary information. This view makes explicit the trade-off studied here: pure predictors are efficient but can be biased under distribution shift, design-based estimators are unbiased but only modestly efficient.

\paragraph{IRT for model evaluation and recalibration.}
Multidimensional IRT is standard in psychometrics and was adopted for LLM evaluation by several of the works above. Recalibrating predicted logits with a fitted slope and intercept goes back to Platt scaling~\citep{platt1999probabilistic}; we apply it per target model, to the IRT logits.
